%% GENERATED by tools/from_docs.py from stepss-docs. Do not edit by hand:
%% edit the page named below in stepss-docs and regenerate.
%% source: stepss-docs/src/content/docs/python/api-reference.md

\chapter{Python API reference}\label{doc:py-api}

\section{\texorpdfstring{\texttt{stepss.cfg}: Test Case Configuration}{stepss.cfg: Test Case Configuration}}\label{py-api:stepss.cfg-test-case-configuration}

The \texttt{stepss.cfg} class defines a simulation scenario: data files, disturbance file, output files, observables, and runtime options.

\begin{center}\rule{0.5\linewidth}{0.5pt}\end{center}

\subsection{Initializing}\label{py-api:initializing}

Create an empty configuration or load one from a previously saved command file:

\begin{Verbatim}
import stepss

case = stepss.cfg()             # empty configuration
case = stepss.cfg("cmd.txt")    # load from command file
\end{Verbatim}

Load multiple cases in a loop:

\begin{Verbatim}
import stepss

list_of_cases = []
for i in range(12):
    list_of_cases.append(stepss.cfg('cmd' + str(i) + '.txt'))
\end{Verbatim}

\subsubsection{\texorpdfstring{\texttt{writeCmdFile(filename=None)}}{writeCmdFile(filename=None)}}\label{py-api:writecmdfilefilenamenone}

Save the current configuration to a command file. Useful for reproducing a run later. When called without an argument, it returns the command-file content as a string instead of writing a file.

\begin{Verbatim}
case.writeCmdFile('cmd.txt')
\end{Verbatim}

Save multiple cases in a loop:

\begin{Verbatim}
for i in range(12):
    list_of_cases[i].writeCmdFile('cmd' + str(i) + '.txt')
\end{Verbatim}

\begin{center}\rule{0.5\linewidth}{0.5pt}\end{center}

\subsection{Data Files}\label{py-api:data-files}

Data files describe the network topology, dynamic models, and solver settings. At least one must be provided.

\subsubsection{\texorpdfstring{\texttt{addData(filename)}}{addData(filename)}}\label{py-api:adddatafilename}

Add a data file to the case.

\begin{Verbatim}
case.addData('dyn_A.dat')
case.addData('settings1.dat')
\end{Verbatim}

\subsubsection{\texorpdfstring{\texttt{delData(filename)}}{delData(filename)}}\label{py-api:deldatafilename}

Remove a specific data file from the case.

\begin{Verbatim}
case.delData('dyn_A.dat')
\end{Verbatim}

\subsubsection{\texorpdfstring{\texttt{getData()}}{getData()}}\label{py-api:getdata}

Return the list of currently registered data files.

\begin{Verbatim}
files = case.getData()
\end{Verbatim}

\subsubsection{\texorpdfstring{\texttt{clearData()}}{clearData()}}\label{py-api:cleardata}

Remove all data files from the case.

\begin{Verbatim}
case.clearData()
\end{Verbatim}

\begin{docsnote}{Caution}

At least one data file must be provided, otherwise the simulator will raise an exception.

\end{docsnote}

\begin{center}\rule{0.5\linewidth}{0.5pt}\end{center}

\subsection{Initialization File}\label{py-api:initialization-file}

Specifies where the simulator writes initialization procedure output.

\subsubsection{\texorpdfstring{\texttt{addInit(filename)}}{addInit(filename)}}\label{py-api:addinitfilename}

\begin{Verbatim}
case.addInit('init.trace')
\end{Verbatim}

\subsubsection{\texorpdfstring{\texttt{getInit()}}{getInit()}}\label{py-api:getinit}

Return the currently registered initialization file path.

\begin{Verbatim}
path = case.getInit()
\end{Verbatim}

\begin{docsnote}{Note}

This is optional. If omitted, the simulator skips writing initialization output.

\end{docsnote}

\begin{center}\rule{0.5\linewidth}{0.5pt}\end{center}

\subsection{Disturbance File}\label{py-api:disturbance-file}

Describes the disturbances to be simulated (generator trips, faults, parameter changes, etc.).

\subsubsection{\texorpdfstring{\texttt{addDst(filename)}}{addDst(filename)}}\label{py-api:adddstfilename}

\begin{Verbatim}
case.addDst('events.dst')
\end{Verbatim}

\subsubsection{\texorpdfstring{\texttt{getDst()}}{getDst()}}\label{py-api:getdst}

Return the currently registered disturbance file path.

\begin{Verbatim}
path = case.getDst()
\end{Verbatim}

\subsubsection{\texorpdfstring{\texttt{clearDst()}}{clearDst()}}\label{py-api:cleardst}

Remove the disturbance file from the case.

\begin{Verbatim}
case.clearDst()
\end{Verbatim}

\begin{docsnote}{Caution}

A disturbance file must be provided, otherwise the simulator will raise an exception.

\end{docsnote}

\begin{center}\rule{0.5\linewidth}{0.5pt}\end{center}

\subsection{Trajectory File}\label{py-api:trajectory-file}

Specifies the file where time-series simulation results (trajectories) are saved for post-processing. This file is used by \texttt{stepss.extractor} to access results after the simulation completes.

\subsubsection{\texorpdfstring{\texttt{addTrj(filename)}}{addTrj(filename)}}\label{py-api:addtrjfilename}

\begin{Verbatim}
case.addTrj('output.trj')
\end{Verbatim}

\subsubsection{\texorpdfstring{\texttt{getTrj()}}{getTrj()}}\label{py-api:gettrj}

Return the currently registered trajectory file path.

\begin{Verbatim}
path = case.getTrj()
\end{Verbatim}

\begin{docsnote}{Note}

This is optional. If omitted, no trajectory file is written and result extraction via \texttt{stepss.extractor} will not be possible.

\end{docsnote}

\begin{center}\rule{0.5\linewidth}{0.5pt}\end{center}

\subsection{Observables File}\label{py-api:observables-file}

Defines which components and quantities are recorded in the trajectory file.

\subsubsection{\texorpdfstring{\texttt{addObs(filename)}}{addObs(filename)}}\label{py-api:addobsfilename}

\begin{Verbatim}
case.addObs('obs.dat')
\end{Verbatim}

\subsubsection{\texorpdfstring{\texttt{getObs()}}{getObs()}}\label{py-api:getobs}

Return the currently registered observables file path.

\begin{Verbatim}
path = case.getObs()
\end{Verbatim}

\begin{docsnote}{Note}

Required when a trajectory file is specified. Defines what data is stored in the trajectory.

\end{docsnote}

\begin{center}\rule{0.5\linewidth}{0.5pt}\end{center}

\subsection{Output/Trace Files}\label{py-api:outputtrace-files}

\subsubsection{\texorpdfstring{\texttt{addOut(filename)}}{addOut(filename)}}\label{py-api:addoutfilename}

Set the main output trace file for simulation progress logging.

\begin{Verbatim}
case.addOut('output.trace')
\end{Verbatim}

\subsubsection{\texorpdfstring{\texttt{getOut()}}{getOut()}}\label{py-api:getout}

Return the currently registered output trace file path.

\begin{Verbatim}
path = case.getOut()
\end{Verbatim}

\subsubsection{\texorpdfstring{\texttt{addCont(filename)}}{addCont(filename)}}\label{py-api:addcontfilename}

Set the continuous trace file. Records Newton solver convergence information at each step. Useful for debugging but can slow down the simulation.

\begin{Verbatim}
case.addCont('cont.trace')
\end{Verbatim}

\subsubsection{\texorpdfstring{\texttt{getCont()}}{getCont()}}\label{py-api:getcont}

Return the currently registered continuous trace file path.

\begin{Verbatim}
path = case.getCont()
\end{Verbatim}

\subsubsection{\texorpdfstring{\texttt{addDisc(filename)}}{addDisc(filename)}}\label{py-api:adddiscfilename}

Set the discrete trace file. Records discrete events: switching actions from disturbance files, discrete controllers, or discrete variables in injector/torque/exciter/two-port models.

\begin{Verbatim}
case.addDisc('disc.trace')
\end{Verbatim}

\subsubsection{\texorpdfstring{\texttt{getDisc()}}{getDisc()}}\label{py-api:getdisc}

Return the currently registered discrete trace file path.

\begin{Verbatim}
path = case.getDisc()
\end{Verbatim}

\subsubsection{\texorpdfstring{\texttt{clearDisc()}}{clearDisc()}}\label{py-api:cleardisc}

Remove the discrete trace file from the case.

\begin{Verbatim}
case.clearDisc()
\end{Verbatim}

\begin{docsnote}{Note}

All output/trace files are optional.

\end{docsnote}

\begin{center}\rule{0.5\linewidth}{0.5pt}\end{center}

\subsection{Runtime Observables}\label{py-api:runtime-observables}

Runtime observables are recorded by the engine while a simulation runs, into a
curve file you can read with \texttt{cur} and plot with \texttt{curplot}. Nothing has to be
installed for this.

\subsubsection{\texorpdfstring{\texttt{addRunObs(obs\_string)}}{addRunObs(obs\_string)}}\label{py-api:addrunobsobs_string}

Add a runtime observable. The following observable types are supported:

\textbf{\texttt{BV\ BUSNAME}}, Voltage magnitude of a bus:

\begin{Verbatim}
case.addRunObs('BV 1041')
\end{Verbatim}

\textbf{\texttt{MS\ MACHINE\_NAME}}, Rotor speed of a synchronous machine:

\begin{Verbatim}
case.addRunObs('MS g1')
\end{Verbatim}

\textbf{\texttt{BPE\ /\ BQE\ /\ BPO\ /\ BQO\ BRANCH\_NAME}}, Active (P) or reactive (Q) power at the origin (O) or extremity (E) of a branch:

\begin{Verbatim}
case.addRunObs('BPO 1041-01')   # active power at origin of branch 1041-01
case.addRunObs('BQO 1041-01')   # reactive power at origin
case.addRunObs('BPE 1041-01')   # active power at extremity
case.addRunObs('BQE 1041-01')   # reactive power at extremity
\end{Verbatim}

\textbf{\texttt{ON\ INJECTOR\_NAME\ OBSERVABLE\_NAME}}, Named observable from an injector model:

\begin{Verbatim}
case.addRunObs('ON WT1a Pw')    # observable Pw from injector WT1a
\end{Verbatim}

\textbf{\texttt{TO\ TWOP\_NAME\ OBSERVABLE\_NAME}}, Named observable from a two-port model:

\begin{Verbatim}
case.addRunObs('TO hvdc1 P1')   # observable P1 from two-port hvdc1
\end{Verbatim}

\textbf{\texttt{RT\ RT}}, Real-time versus simulated-time plot (useful to gauge simulation speed):

\begin{Verbatim}
case.addRunObs('RT RT')
\end{Verbatim}

\subsubsection{\texorpdfstring{\texttt{clearRunObs()}}{clearRunObs()}}\label{py-api:clearrunobs}

Remove all runtime observables.

\begin{Verbatim}
case.clearRunObs()
\end{Verbatim}

\begin{docsnote}{Note}

Nothing needs to be installed for run-time observables.

\end{docsnote}

\begin{center}\rule{0.5\linewidth}{0.5pt}\end{center}

\section{\texorpdfstring{\texttt{stepss.sim}: Simulation Control}{stepss.sim: Simulation Control}}\label{py-api:stepss.sim-simulation-control}

The \texttt{stepss.sim} class runs simulations. It wraps the RAMSES dynamic library and supports start/pause/continue, runtime queries, and disturbance injection.

\begin{center}\rule{0.5\linewidth}{0.5pt}\end{center}

\subsection{Initializing}\label{py-api:initializing-1}

\begin{Verbatim}
import stepss

ram = stepss.sim()                        # use bundled RAMSES libraries
ram = stepss.sim(custLibDir='/path/to/')  # use custom library directory
\end{Verbatim}

{\def\LTcaptype{none} % do not increment counter
\begin{small}\begin{longtable}[]{@{}
  >{\raggedright\arraybackslash}p{(\linewidth - 6\tabcolsep) * \real{0.2482}}
  >{\raggedright\arraybackslash}p{(\linewidth - 6\tabcolsep) * \real{0.2206}}
  >{\raggedright\arraybackslash}p{(\linewidth - 6\tabcolsep) * \real{0.5312}}
@{}}
\toprule\noalign{}
\begin{minipage}[b]{\linewidth}\raggedright
Parameter
\end{minipage} & \begin{minipage}[b]{\linewidth}\raggedright
Type
\end{minipage} & \begin{minipage}[b]{\linewidth}\raggedright
Description
\end{minipage} \\
\midrule\noalign{}
\endhead
\bottomrule\noalign{}
\endlastfoot
\texttt{custLibDir} & \texttt{str} or \texttt{None} & Custom path to the RAMSES library directory. Default: use bundled libraries. \\
\end{longtable}\end{small}
}

\begin{center}\rule{0.5\linewidth}{0.5pt}\end{center}

\subsection{Running Simulations}\label{py-api:running-simulations}

A properly configured \texttt{stepss.cfg} test case is required before running a simulation.

\subsubsection{\texorpdfstring{\texttt{execSim(case)}: run to completion}{execSim(case): run to completion}}\label{py-api:execsimcase-run-to-completion}

\begin{Verbatim}
ram.execSim(case)
\end{Verbatim}

\subsubsection{\texorpdfstring{\texttt{execSim(case,\ t)}: start and pause at time t}{execSim(case, t): start and pause at time t}}\label{py-api:execsimcase-t-start-and-pause-at-time-t}

Start the simulation and pause at a specific time (in seconds):

\begin{Verbatim}
ram.execSim(case, 10.0)    # start and pause at t = 10 s
\end{Verbatim}

\subsubsection{\texorpdfstring{\texttt{contSim(t)}: continue to time t}{contSim(t): continue to time t}}\label{py-api:contsimt-continue-to-time-t}

Resume a paused simulation until a specified time:

\begin{Verbatim}
ram.contSim(20.0)                     # resume until t = 20 s
ram.contSim(ram.getSimTime() + 60.0)  # advance by 60 s from current time
ram.contSim(ram.getInfTime())         # run to the end of the time horizon
\end{Verbatim}

\subsubsection{\texorpdfstring{\texttt{endSim()}: terminate early}{endSim(): terminate early}}\label{py-api:endsim-terminate-early}

Terminate the simulation before reaching the time horizon:

\begin{Verbatim}
ram.endSim()
\end{Verbatim}

\subsubsection{\texorpdfstring{\texttt{pauseSim(t\_pause)}: schedule a pause}{pauseSim(t\_pause): schedule a pause}}\label{py-api:pausesimt_pause-schedule-a-pause}

Schedule a pause at a given simulated time; it takes effect on the next \texttt{execSim()} or \texttt{contSim()} call:

\begin{Verbatim}
ram.pauseSim(10.0)
ram.execSim(case)   # will pause at t = 10 s
\end{Verbatim}

\subsubsection{\texorpdfstring{\texttt{getEndSim()}}{getEndSim()}}\label{py-api:getendsim}

Return \texttt{0} while the simulation is still running, \texttt{1} once it has ended:

\begin{Verbatim}
if ram.getEndSim():
    print('simulation finished')
\end{Verbatim}

\subsubsection{\texorpdfstring{\texttt{getLastErr()}}{getLastErr()}}\label{py-api:getlasterr}

Return the last error message issued by RAMSES as a string. Useful after catching a \texttt{RAMSESError}:

\begin{Verbatim}
msg = ram.getLastErr()
\end{Verbatim}

\begin{center}\rule{0.5\linewidth}{0.5pt}\end{center}

\subsection{Querying State}\label{py-api:querying-state}

When the simulation is paused, the following methods query the current system state.

\subsubsection{\texorpdfstring{\texttt{getSimTime()}}{getSimTime()}}\label{py-api:getsimtime}

Return the current simulation time in seconds.

\begin{Verbatim}
t = ram.getSimTime()
\end{Verbatim}

\subsubsection{\texorpdfstring{\texttt{getInfTime()}}{getInfTime()}}\label{py-api:getinftime}

Return the value used as ``infinity'' (i.e., the end of the simulation time horizon). Pass this to \texttt{contSim()} to run to completion.

\begin{Verbatim}
t_inf = ram.getInfTime()
ram.contSim(ram.getInfTime())
\end{Verbatim}

\subsubsection{\texorpdfstring{\texttt{getAllCompNames(type)}}{getAllCompNames(type)}}\label{py-api:getallcompnamestype}

Return a list of all component names of the given type.

\begin{Verbatim}
buses    = ram.getAllCompNames('BUS')     # list of all bus names
gens     = ram.getAllCompNames('SYNC')    # list of all generator names
injs     = ram.getAllCompNames('INJ')     # list of all injector names
dctls    = ram.getAllCompNames('DCTL')    # list of all discrete controller names
branches = ram.getAllCompNames('BRANCH')  # list of all branch names
twops    = ram.getAllCompNames('TWOP')    # list of all two-port names
shunts   = ram.getAllCompNames('SHUNT')   # list of all shunt names
loads    = ram.getAllCompNames('LOAD')    # list of all load names
\end{Verbatim}

Supported component types: \texttt{BUS}, \texttt{SYNC}, \texttt{INJ}, \texttt{DCTL}, \texttt{BRANCH}, \texttt{TWOP}, \texttt{SHUNT}, \texttt{LOAD}.

\subsubsection{\texorpdfstring{\texttt{getCompName(comp\_type,\ num)}}{getCompName(comp\_type, num)}}\label{py-api:getcompnamecomp_type-num}

Return the name of the \texttt{num}-th component (1-based) of the given type. Accepts the same component types as \texttt{getAllCompNames()}.

\begin{Verbatim}
first_bus = ram.getCompName('BUS', 1)   # e.g. 'B1'
\end{Verbatim}

\subsubsection{\texorpdfstring{\texttt{getBusVolt(names)}}{getBusVolt(names)}}\label{py-api:getbusvoltnames}

Return voltage magnitudes (in pu) for a list of bus names.

\begin{Verbatim}
ram.execSim(case, 10.0)
bus_names = ['g1', 'g2', '4032']
voltages = ram.getBusVolt(bus_names)
\end{Verbatim}

\subsubsection{\texorpdfstring{\texttt{getBusPha(names)}}{getBusPha(names)}}\label{py-api:getbusphanames}

Return voltage phase angles (in degrees) for a list of bus names.

\begin{Verbatim}
phases = ram.getBusPha(bus_names)
\end{Verbatim}

\subsubsection{\texorpdfstring{\texttt{getBranchPow(names)}}{getBranchPow(names)}}\label{py-api:getbranchpownames}

Return power flows for a list of branch names. Each entry is \texttt{{[}P\_from,\ Q\_from,\ P\_to,\ Q\_to{]}} in MW and Mvar.

\begin{Verbatim}
powers = ram.getBranchPow(['1041-01'])
# powers[0] == [P_from, Q_from, P_to, Q_to]
\end{Verbatim}

\subsubsection{\texorpdfstring{\texttt{getBranchCur(names)}}{getBranchCur(names)}}\label{py-api:getbranchcurnames}

Return branch currents for a list of branch names. Each entry contains the x--y current components at the origin and extremity: \texttt{{[}ix\_orig,\ iy\_orig,\ ix\_extr,\ iy\_extr{]}}.

\begin{Verbatim}
currents = ram.getBranchCur(['1011-1013', '1012-1014'])
\end{Verbatim}

\subsubsection{\texorpdfstring{\texttt{getObs(comp\_types,\ comp\_names,\ obs\_names)}}{getObs(comp\_types, comp\_names, obs\_names)}}\label{py-api:getobscomp_types-comp_names-obs_names}

Get the current value of named observables for a list of components. Lists must be the same length.

\begin{Verbatim}
comp_type = ['INJ', 'EXC', 'TOR']
comp_name = ['L_11', 'g2',  'g3']
obs_name  = ['P',    'vf',  'Pm']
obs = ram.getObs(comp_type, comp_name, obs_name)
\end{Verbatim}

Supported model types: \texttt{EXC} (exciter), \texttt{TOR} (governor), \texttt{INJ} (injector), \texttt{TWOP} (two-port), \texttt{DCTL} (discrete controller), \texttt{SYN} (synchronous generator).

\subsubsection{\texorpdfstring{\texttt{getPrm(comp\_types,\ comp\_names,\ prm\_names)}}{getPrm(comp\_types, comp\_names, prm\_names)}}\label{py-api:getprmcomp_types-comp_names-prm_names}

Get parameter values for a list of components. Lists must be the same length.

\begin{Verbatim}
comp_type = ['EXC', 'EXC']
comp_name = ['g1',  'g2']
prm_name  = ['V0',  'KPSS']
prms = ram.getPrm(comp_type, comp_name, prm_name)
\end{Verbatim}

\subsubsection{\texorpdfstring{\texttt{getPrmNames(comp\_types,\ comp\_names)}}{getPrmNames(comp\_types, comp\_names)}}\label{py-api:getprmnamescomp_types-comp_names}

List the parameter names of one or more model instances. Accepts single strings or equal-length lists; component types are \texttt{EXC}, \texttt{TOR}, \texttt{INJ}, \texttt{DCTL}, \texttt{TWOP}.

\begin{Verbatim}
names = ram.getPrmNames('EXC', 'g1')   # list of parameter names of g1's exciter
\end{Verbatim}

\begin{center}\rule{0.5\linewidth}{0.5pt}\end{center}

\subsection{Runtime Observable Recording}\label{py-api:runtime-observable-recording}

Observables can also be selected programmatically after the simulation starts, instead of via an observables file. Call the three methods in order, after \texttt{execSim(case,\ 0.0)}:

\subsubsection{\texorpdfstring{\texttt{initObserv(traj\_filenm)}}{initObserv(traj\_filenm)}}\label{py-api:initobservtraj_filenm}

Initialize the runtime observable recording system and set the output trajectory file.

\subsubsection{\texorpdfstring{\texttt{addObserv(string)}}{addObserv(string)}}\label{py-api:addobservstring}

Register one observable selector in RAMSES format (e.g.~\texttt{\textquotesingle{}BUS\ *\textquotesingle{}}, \texttt{\textquotesingle{}SYNC\ g1\textquotesingle{}}). Call once per selector.

\subsubsection{\texorpdfstring{\texttt{finalObserv()}}{finalObserv()}}\label{py-api:finalobserv}

Finalize the selection, allocate the recording buffers, and write the trajectory file header.

\begin{Verbatim}
ram.execSim(case, 0.0)
ram.initObserv('obs.trj')
ram.addObserv('BUS *')      # record all bus voltages
ram.addObserv('SYNC g1')    # record machine g1
ram.finalObserv()
ram.contSim(ram.getInfTime())
\end{Verbatim}

\begin{center}\rule{0.5\linewidth}{0.5pt}\end{center}

\subsection{Subsystems}\label{py-api:subsystems}

Subsystems select a group of buses for aggregate queries.

\subsubsection{\texorpdfstring{\texttt{defineSS(ssID,\ filter1,\ filter2,\ filter3)}}{defineSS(ssID, filter1, filter2, filter3)}}\label{py-api:definessssid-filter1-filter2-filter3}

Define subsystem \texttt{ssID} as the intersection of three filters: voltage levels, zones, and bus names. Each filter is a list of strings; an empty list deactivates that filter.

\begin{Verbatim}
ram.defineSS(1, ['735'], [], [])   # subsystem 1 = all 735 kV buses
\end{Verbatim}

\subsubsection{\texorpdfstring{\texttt{getSS(ssID)}}{getSS(ssID)}}\label{py-api:getssssid}

Return the list of buses belonging to a subsystem.

\begin{Verbatim}
buses = ram.getSS(1)
\end{Verbatim}

\subsubsection{\texorpdfstring{\texttt{getTrfoSS(ssID,\ location,\ in\_service,\ rettype)}}{getTrfoSS(ssID, location, in\_service, rettype)}}\label{py-api:gettrfossssid-location-in_service-rettype}

Return transformer information for a subsystem. \texttt{location}: \texttt{1} = both ends inside the subsystem, \texttt{2} = tie transformers, \texttt{3} = both. \texttt{in\_service}: \texttt{1} = in-service only, \texttt{2} = all. \texttt{rettype} selects the returned quantity: \texttt{NAME}, \texttt{From}, \texttt{To}, \texttt{Status}, \texttt{Tap}, \texttt{Currentf}, \texttt{Currentt}, \texttt{Pf}, \texttt{Qf}, \texttt{Pt}, \texttt{Qt}.

\begin{Verbatim}
status = ram.getTrfoSS(1, 3, 2, 'Status')
\end{Verbatim}

\begin{center}\rule{0.5\linewidth}{0.5pt}\end{center}

\subsection{Runtime Disturbances}\label{py-api:runtime-disturbances}

Disturbances can be added dynamically while the simulation is paused, enabling interactive scenario analysis.

\subsubsection{\texorpdfstring{\texttt{addDisturb(time,\ description)}}{addDisturb(time, description)}}\label{py-api:adddisturbtime-description}

Schedule a disturbance to occur at a given simulation time.

The disturbance description string follows the same syntax as the disturbance file format. See \hyperref[ch-disturb]{Disturbances} for the complete reference.

\begin{Verbatim}
ram.execSim(case, 80.0)

# Trip generator g7 at t = 100 s
ram.addDisturb(100.0, 'BREAKER SYNC_MACH g7 0')

# Apply a 3-phase fault at bus 4032 and clear it 100 ms later
ram.addDisturb(100.0, 'FAULT BUS 4032 0. 0.')
ram.addDisturb(100.1, 'CLEAR BUS 4032')

# Step change in an LTC setpoint
ram.addDisturb(100.0, 'CHGPRM DCTL 1-1041 Vsetpt -0.05 0')

ram.contSim(ram.getInfTime())
\end{Verbatim}

\begin{center}\rule{0.5\linewidth}{0.5pt}\end{center}

\subsection{Jacobian Export}\label{py-api:jacobian-export}

\subsubsection{\texorpdfstring{\texttt{getJac()}}{getJac()}}\label{py-api:getjac}

Export the system Jacobian in descriptor form at the current pause point. Returns a tuple \texttt{(A,\ E)} of \texttt{scipy.sparse.csc\_matrix} objects, where \texttt{A} holds the numerical Jacobian values and \texttt{E} is the structural incidence matrix of the descriptor system. Two intermediate text files are also written to the working directory:

{\def\LTcaptype{none} % do not increment counter
\begin{small}\begin{longtable}[]{@{}ll@{}}
\toprule\noalign{}
File & Contents \\
\midrule\noalign{}
\endhead
\bottomrule\noalign{}
\endlastfoot
\texttt{py\_val.dat} & Jacobian values (parsed into \texttt{A}) \\
\texttt{py\_eqs.dat} & Structural incidence matrix (parsed into \texttt{E}) \\
\end{longtable}\end{small}
}

\begin{Verbatim}
ram.execSim(case, 10.0)
A, E = ram.getJac()
\end{Verbatim}

Use this when you want to drive your own solver, for instance the sparse
shift-invert methods in \texttt{scipy.sparse.linalg} that large systems need. To have
RAMSES do the analysis instead, schedule an \texttt{EIG} disturbance; see
\hyperref[doc:eigenanalysis]{Eigenanalysis}.

\begin{docsnote}{Note}

Set \texttt{\$OMEGA\_REF\ SYN\ ;} in the solver settings data file when exporting the Jacobian for eigenanalysis.

\end{docsnote}

\begin{center}\rule{0.5\linewidth}{0.5pt}\end{center}

\section{\texorpdfstring{\texttt{stepss.extractor}: Result Extraction}{stepss.extractor: Result Extraction}}\label{py-api:stepss.extractor-result-extraction}

The \texttt{stepss.extractor} class extracts and visualises time-series results from a trajectory file produced during simulation.

\begin{center}\rule{0.5\linewidth}{0.5pt}\end{center}

\subsection{Initializing}\label{py-api:initializing-2}

Pass the trajectory file path to the extractor:

\begin{Verbatim}
import stepss

case = stepss.cfg('cmd.txt')
# ... run simulation ...
ext = stepss.extractor(case.getTrj())
\end{Verbatim}

Or provide the file path directly:

\begin{Verbatim}
ext = stepss.extractor('output.trj')
\end{Verbatim}

\begin{center}\rule{0.5\linewidth}{0.5pt}\end{center}

\subsection{Curve Objects}\label{py-api:curve-objects}

All extraction methods return objects whose attributes are \textbf{curve objects} (\texttt{stepss.cur} named tuples). Every curve object has:

{\def\LTcaptype{none} % do not increment counter
\begin{small}\begin{longtable}[]{@{}lll@{}}
\toprule\noalign{}
Attribute & Type & Description \\
\midrule\noalign{}
\endhead
\bottomrule\noalign{}
\endlastfoot
\texttt{time} & \texttt{numpy.ndarray} & Time values in seconds \\
\texttt{value} & \texttt{numpy.ndarray} & Observable values \\
\texttt{msg} & \texttt{str} & Description string (used as plot legend label) \\
\end{longtable}\end{small}
}

\subsubsection{\texorpdfstring{\texttt{.plot()}}{.plot()}}\label{py-api:plot}

Display the curve using Matplotlib:

\begin{Verbatim}
bus = ext.getBus('4044')
bus.mag.plot()
\end{Verbatim}

\begin{center}\rule{0.5\linewidth}{0.5pt}\end{center}

\subsection{Extraction Methods}\label{py-api:extraction-methods}

\subsubsection{\texorpdfstring{\texttt{getBus(name)}}{getBus(name)}}\label{py-api:getbusname}

Retrieve voltage time series for a bus. Returns an object with:

{\def\LTcaptype{none} % do not increment counter
\begin{small}\begin{longtable}[]{@{}ll@{}}
\toprule\noalign{}
Attribute & Description \\
\midrule\noalign{}
\endhead
\bottomrule\noalign{}
\endlastfoot
\texttt{.mag} & Voltage magnitude (pu) \\
\texttt{.pha} & Voltage phase angle (deg) \\
\end{longtable}\end{small}
}

\begin{Verbatim}
bus = ext.getBus('4044')
bus.mag.plot()   # voltage magnitude (pu)
bus.pha.plot()   # voltage phase angle (deg)
\end{Verbatim}

\begin{center}\rule{0.5\linewidth}{0.5pt}\end{center}

\subsubsection{\texorpdfstring{\texttt{getSync(name)}}{getSync(name)}}\label{py-api:getsyncname}

Retrieve the full set of synchronous machine observables. Returns an object with:

{\def\LTcaptype{none} % do not increment counter
\begin{small}\begin{longtable}[]{@{}ll@{}}
\toprule\noalign{}
Attribute & Description \\
\midrule\noalign{}
\endhead
\bottomrule\noalign{}
\endlastfoot
\texttt{.P} & Active power (MW) \\
\texttt{.Q} & Reactive power (Mvar) \\
\texttt{.S} & Rotor speed (pu; 1 = nominal) \\
\texttt{.A} & Rotor angle w.r.t. COI (deg) \\
\texttt{.FV} & Field voltage (pu) \\
\texttt{.FC} & Field current (pu) \\
\texttt{.T} & Mechanical torque (pu) \\
\texttt{.ET} & Electromagnetic torque (pu) \\
\texttt{.FW} & Field winding flux \\
\texttt{.DD} & d1 damper flux \\
\texttt{.QD} & q1 damper flux \\
\texttt{.QW} & q2 winding flux \\
\texttt{.SC} & COI speed \textbf{deviation} (pu; 0 = nominal) \\
\end{longtable}\end{small}
}

\begin{docsnote}{Note}

The two speed observables use different conventions: \texttt{.S} is the whole rotor
speed (1 = nominal), while \texttt{.SC} is the deviation of the centre-of-inertia
(COI) speed from nominal (0 = nominal). The system frequency in Hz is
\texttt{fnom\ *\ (1\ +\ SC)} or, per machine, \texttt{fnom\ *\ S}.

\end{docsnote}

\begin{Verbatim}
gen = ext.getSync('g1')
gen.P.plot()    # active power (MW)
gen.Q.plot()    # reactive power (Mvar)
gen.S.plot()    # rotor speed (pu)
gen.A.plot()    # rotor angle w.r.t. COI (deg)
gen.FV.plot()   # field voltage (pu)
gen.FC.plot()   # field current (pu)
gen.T.plot()    # mechanical torque (pu)
gen.ET.plot()   # electromagnetic torque (pu)
\end{Verbatim}

\begin{center}\rule{0.5\linewidth}{0.5pt}\end{center}

\subsubsection{\texorpdfstring{\texttt{getExc(name)}}{getExc(name)}}\label{py-api:getexcname}

Retrieve exciter observables. Available observables depend on the exciter model.

{\def\LTcaptype{none} % do not increment counter
\begin{small}\begin{longtable}[]{@{}ll@{}}
\toprule\noalign{}
Attribute & Description \\
\midrule\noalign{}
\endhead
\bottomrule\noalign{}
\endlastfoot
\texttt{.obsdict} & \texttt{dict} mapping observable name \ensuremath{\rightarrow} description \\
\emph{(model-dependent)} & Access by observable name, e.g.~\texttt{.vf} \\
\end{longtable}\end{small}
}

\begin{Verbatim}
exc = ext.getExc('g1')
print(exc.obsdict)   # list available observables for this model
exc.vf.plot()        # field voltage (model-dependent name)
\end{Verbatim}

\begin{center}\rule{0.5\linewidth}{0.5pt}\end{center}

\subsubsection{\texorpdfstring{\texttt{getTor(name)}}{getTor(name)}}\label{py-api:gettorname}

Retrieve governor/torque model observables. Available observables depend on the governor model.

{\def\LTcaptype{none} % do not increment counter
\begin{small}\begin{longtable}[]{@{}ll@{}}
\toprule\noalign{}
Attribute & Description \\
\midrule\noalign{}
\endhead
\bottomrule\noalign{}
\endlastfoot
\texttt{.obsdict} & \texttt{dict} mapping observable name \ensuremath{\rightarrow} description \\
\emph{(model-dependent)} & Access by observable name, e.g.~\texttt{.Pm} \\
\end{longtable}\end{small}
}

\begin{Verbatim}
gov = ext.getTor('g1')
print(gov.obsdict)   # list available observables for this model
gov.Pm.plot()        # mechanical power (pu)
\end{Verbatim}

\begin{center}\rule{0.5\linewidth}{0.5pt}\end{center}

\subsubsection{\texorpdfstring{\texttt{getInj(name)}}{getInj(name)}}\label{py-api:getinjname}

Retrieve injector observables. Injectors include renewable energy sources (wind, PV, BESS), loads, and other single-bus components.

{\def\LTcaptype{none} % do not increment counter
\begin{small}\begin{longtable}[]{@{}ll@{}}
\toprule\noalign{}
Attribute & Description \\
\midrule\noalign{}
\endhead
\bottomrule\noalign{}
\endlastfoot
\texttt{.obsdict} & \texttt{dict} mapping observable name \ensuremath{\rightarrow} description \\
\emph{(model-dependent)} & Access by observable name, e.g.~\texttt{.Pw} \\
\end{longtable}\end{small}
}

\begin{Verbatim}
inj = ext.getInj('WT1a')
print(inj.obsdict)   # list available observables
inj.Pw.plot()        # wind power (model-dependent name)
\end{Verbatim}

\begin{center}\rule{0.5\linewidth}{0.5pt}\end{center}

\subsubsection{\texorpdfstring{\texttt{getTwop(name)}}{getTwop(name)}}\label{py-api:gettwopname}

Retrieve two-port model observables. Two-port models include HVDC links (LCC and VSC), SVCs, and DC systems.

{\def\LTcaptype{none} % do not increment counter
\begin{small}\begin{longtable}[]{@{}ll@{}}
\toprule\noalign{}
Attribute & Description \\
\midrule\noalign{}
\endhead
\bottomrule\noalign{}
\endlastfoot
\texttt{.obsdict} & \texttt{dict} mapping observable name \ensuremath{\rightarrow} description \\
\emph{(model-dependent)} & Access by observable name, e.g.~\texttt{.P1}, \texttt{.P2} \\
\end{longtable}\end{small}
}

\begin{Verbatim}
twop = ext.getTwop('hvdc1')
print(twop.obsdict)  # list available observables
twop.P1.plot()       # active power at terminal 1
twop.P2.plot()       # active power at terminal 2
\end{Verbatim}

\begin{center}\rule{0.5\linewidth}{0.5pt}\end{center}

\subsubsection{\texorpdfstring{\texttt{getDctl(name)}}{getDctl(name)}}\label{py-api:getdctlname}

Retrieve discrete controller observables. Discrete controllers include LTC transformers, under-voltage load shedding, phase shifters, etc.

{\def\LTcaptype{none} % do not increment counter
\begin{small}\begin{longtable}[]{@{}ll@{}}
\toprule\noalign{}
Attribute & Description \\
\midrule\noalign{}
\endhead
\bottomrule\noalign{}
\endlastfoot
\texttt{.obsdict} & \texttt{dict} mapping observable name \ensuremath{\rightarrow} description \\
\end{longtable}\end{small}
}

\begin{Verbatim}
dctl = ext.getDctl('1-1041')
print(dctl.obsdict)  # list available observables
\end{Verbatim}

\begin{center}\rule{0.5\linewidth}{0.5pt}\end{center}

\subsubsection{\texorpdfstring{\texttt{getBranch(name)}}{getBranch(name)}}\label{py-api:getbranchname}

Retrieve branch (line/transformer) power flow time series.

{\def\LTcaptype{none} % do not increment counter
\begin{small}\begin{longtable}[]{@{}ll@{}}
\toprule\noalign{}
Attribute & Description \\
\midrule\noalign{}
\endhead
\bottomrule\noalign{}
\endlastfoot
\texttt{.PF} & Active power at FROM end (MW) \\
\texttt{.QF} & Reactive power at FROM end (Mvar) \\
\texttt{.PT} & Active power at TO end (MW) \\
\texttt{.QT} & Reactive power at TO end (Mvar) \\
\texttt{.RM} & Transformer ratio magnitude \\
\texttt{.RA} & Transformer phase angle (deg) \\
\end{longtable}\end{small}
}

\begin{Verbatim}
branch = ext.getBranch('1041-01')
branch.PF.plot()   # active power at FROM end (MW)
branch.QF.plot()   # reactive power at FROM end (Mvar)
branch.PT.plot()   # active power at TO end (MW)
branch.QT.plot()   # reactive power at TO end (Mvar)
branch.RM.plot()   # transformer ratio magnitude
branch.RA.plot()   # transformer phase angle (deg)
\end{Verbatim}

\begin{center}\rule{0.5\linewidth}{0.5pt}\end{center}

\subsubsection{\texorpdfstring{\texttt{getShunt(name)}}{getShunt(name)}}\label{py-api:getshuntname}

Retrieve shunt compensation time series.

{\def\LTcaptype{none} % do not increment counter
\begin{small}\begin{longtable}[]{@{}ll@{}}
\toprule\noalign{}
Attribute & Description \\
\midrule\noalign{}
\endhead
\bottomrule\noalign{}
\endlastfoot
\texttt{.Q} & Reactive power produced (Mvar) \\
\end{longtable}\end{small}
}

\begin{Verbatim}
ext.getShunt('sh1').Q.plot()
\end{Verbatim}

\begin{center}\rule{0.5\linewidth}{0.5pt}\end{center}

\subsubsection{\texorpdfstring{\texttt{getLoad(name)}}{getLoad(name)}}\label{py-api:getloadname}

Retrieve load time series.

{\def\LTcaptype{none} % do not increment counter
\begin{small}\begin{longtable}[]{@{}ll@{}}
\toprule\noalign{}
Attribute & Description \\
\midrule\noalign{}
\endhead
\bottomrule\noalign{}
\endlastfoot
\texttt{.P} & Active power consumed (MW) \\
\texttt{.Q} & Reactive power consumed (Mvar) \\
\end{longtable}\end{small}
}

\begin{Verbatim}
ext.getLoad('L_1').P.plot()
\end{Verbatim}

\begin{center}\rule{0.5\linewidth}{0.5pt}\end{center}

\subsection{Multi-Curve Plotting}\label{py-api:multi-curve-plotting}

\subsubsection{\texorpdfstring{\texttt{stepss.curplot(curves)}}{stepss.curplot(curves)}}\label{py-api:stepss.curplotcurves}

Display multiple curve objects on the same axes. Each curve's \texttt{msg} field is used as the legend label.

\begin{Verbatim}
import stepss

ext = stepss.extractor(case.getTrj())

curves = [
    ext.getSync('g1').S,
    ext.getSync('g2').S,
    ext.getSync('g3').S,
]
stepss.curplot(curves)
\end{Verbatim}

\begin{center}\rule{0.5\linewidth}{0.5pt}\end{center}

\section{\texorpdfstring{\texttt{stepss.monitor}: Live Plotting}{stepss.monitor: Live Plotting}}\label{py-api:stepss.monitor-live-plotting}

\texttt{monitor} plots chosen quantities while a simulation runs. It steps the engine
forward in slices, reads those quantities at every pause, and redraws a stacked
figure: one panel per observable, all sharing the time axis.

The engine is polled in process, so nothing is read from disk and nothing has to
be installed beyond matplotlib, which pip installs with the package. Runtime
observables, above, are a separate mechanism in which the engine writes a curve
file of its own; the two are independent and can be used together.

\begin{Verbatim}
import stepss

ram = stepss.sim()
case = stepss.cfg('cmd.txt')
ram.execSim(case, 0.0)                    # initialise, paused at t = 0

mon = stepss.monitor(ram, ['BV 4044', 'MS g6', 'RT RT'], title='Nordic')
curves = mon.run(step=0.5)                # run to the end of the scenario
\end{Verbatim}

The simulation must already be initialised and paused, which \texttt{execSim} does when
given a \texttt{pause} time.

\subsection{\texorpdfstring{\texttt{monitor(ram,\ observables,\ title=None,\ refresh=0.2,\ show=True)}}{monitor(ram, observables, title=None, refresh=0.2, show=True)}}\label{py-api:monitorram-observables-titlenone-refresh0.2-showtrue}

{\def\LTcaptype{none} % do not increment counter
\begin{small}\begin{longtable}[]{@{}
  >{\raggedright\arraybackslash}p{(\linewidth - 4\tabcolsep) * \real{0.2506}}
  >{\raggedright\arraybackslash}p{(\linewidth - 4\tabcolsep) * \real{0.7494}}
@{}}
\toprule\noalign{}
\begin{minipage}[b]{\linewidth}\raggedright
Argument
\end{minipage} & \begin{minipage}[b]{\linewidth}\raggedright
Description
\end{minipage} \\
\midrule\noalign{}
\endhead
\bottomrule\noalign{}
\endlastfoot
\texttt{ram} & The \texttt{stepss.sim} driving the run. \\
\texttt{observables} & What to plot: descriptor strings, \texttt{(label,\ callable)} pairs, or bare callables. A single observable need not be wrapped in a list. \\
\texttt{title} & Figure title. \\
\texttt{refresh} & Minimum wall-clock seconds between redraws. Samples are never skipped, only draws; \texttt{0} redraws at every sample. \\
\texttt{show} & \texttt{False} collects the samples and builds no figure. \\
\end{longtable}\end{small}
}

\subsection{Observable descriptors}\label{py-api:observable-descriptors}

The vocabulary is the one \texttt{addRunObs} uses, plus \texttt{BA} and the generic \texttt{OBS}:

{\def\LTcaptype{none} % do not increment counter
\begin{small}\begin{longtable}[]{@{}
  >{\raggedright\arraybackslash}p{(\linewidth - 6\tabcolsep) * \real{0.3316}}
  >{\raggedright\arraybackslash}p{(\linewidth - 6\tabcolsep) * \real{0.4600}}
  >{\raggedright\arraybackslash}p{(\linewidth - 6\tabcolsep) * \real{0.2085}}
@{}}
\toprule\noalign{}
\begin{minipage}[b]{\linewidth}\raggedright
Descriptor
\end{minipage} & \begin{minipage}[b]{\linewidth}\raggedright
Quantity
\end{minipage} & \begin{minipage}[b]{\linewidth}\raggedright
Unit
\end{minipage} \\
\midrule\noalign{}
\endhead
\bottomrule\noalign{}
\endlastfoot
\texttt{BV\ BUSNAME} & Voltage magnitude of a bus & pu \\
\texttt{BA\ BUSNAME} & Voltage phase of a bus & deg \\
\texttt{MS\ MACHINE\_NAME} & Rotor speed of a synchronous machine & pu \\
\texttt{BPO\ /\ BQO\ /\ BPE\ /\ BQE\ BRANCH\_NAME} & Active or reactive power at the origin or extremity of a branch & MW, Mvar \\
\texttt{ON\ INJECTOR\_NAME\ OBSERVABLE\_NAME} & Named observable of an injector model & set by the model \\
\texttt{TO\ TWOP\_NAME\ OBSERVABLE\_NAME} & Named observable of a two-port model & set by the model \\
\texttt{OBS\ TYPE\ NAME\ OBSERVABLE\_NAME} & Named observable of any component type \texttt{getObs} accepts & set by the model \\
\texttt{RT\ RT} & Elapsed wall-clock time, to gauge simulation speed & s \\
\end{longtable}\end{small}
}

Anything the descriptors do not cover is a callable, taking the simulator and
returning one number:

\begin{Verbatim}
mon = stepss.monitor(ram, [
    'BV 4044',
    ('4044 reactive reserve (Mvar)', lambda r: r.getObs('SYN', 'g6', 'Q')[0]),
])
\end{Verbatim}

\subsection{\texorpdfstring{\texttt{run(step=1.0,\ until=None)}}{run(step=1.0, until=None)}}\label{py-api:runstep1.0-untilnone}

Advance the simulation, sampling and redrawing at each pause, and return one
\texttt{cur} per observable. Returns when the engine reports the end of the scenario,
when \texttt{until} is reached, or when a slice fails to advance the simulated time.

\begin{Verbatim}
curves = mon.run(step=0.1, until=30.0)    # 30 s of simulated time, sampled every 0.1 s
\end{Verbatim}

\texttt{step} is simulated seconds per slice, so it sets both the sampling resolution
and the redraw cadence. The engine pauses at the first internal step at or after
the requested time, so the samples land on or just after the multiples of \texttt{step}.

If the engine raises, the samples taken up to that point stay available from
\texttt{curves()}:

\begin{Verbatim}
from stepss.globals import RAMSESError

try:
    mon.run(step=1.0)
except RAMSESError:
    stepss.curplot(mon.curves())          # everything up to the failure
\end{Verbatim}

\subsection{\texorpdfstring{\texttt{curves()}}{curves()}}\label{py-api:curves}

Return everything sampled so far, one \texttt{cur} per observable, in the order given
to the constructor. These are the same objects \texttt{extractor} produces, so
\texttt{curplot} and the rest of the post-processing accept them unchanged.

\subsection{\texorpdfstring{\texttt{sample()}}{sample()}}\label{py-api:sample}

Read every observable once, at the current simulated time. Call it directly when
driving the simulation yourself:

\begin{Verbatim}
mon = stepss.monitor(ram, ['BV 4044'])
for target in [10.0, 20.0, 30.0]:
    ram.contSim(target)
    ram.addDisturb(target + 1.0, 'CHGPRM DCTL 1-1041 Vsetpt -0.01 0')
    mon.sample()
    mon.refresh()
\end{Verbatim}

\subsection{\texorpdfstring{\texttt{refresh(force=False)}}{refresh(force=False)}}\label{py-api:refreshforcefalse}

Push the samples into the figure and redraw it. A draw falling inside the
\texttt{refresh} interval is skipped unless \texttt{force} is set.

\subsection{\texorpdfstring{\texttt{savefig(fname,\ **kwargs)}}{savefig(fname, **kwargs)}}\label{py-api:savefigfname-kwargs}

Redraw and write the figure to a file, forwarding to
\texttt{matplotlib.figure.Figure.savefig}.

\subsection{\texorpdfstring{\texttt{close()}}{close()}}\label{py-api:close}

Close the figure. The samples are unaffected: \texttt{curves()} keeps working.

\begin{docsnote}{Note}

On a file-only matplotlib backend such as \texttt{Agg}, and in Jupyter's inline mode,
there is no window to animate: the monitor collects the samples and builds the
figure, and \texttt{savefig} writes the same chart a window would have shown. Closing
the window mid-run stops the redraws and leaves the run going.

\end{docsnote}

\begin{center}\rule{0.5\linewidth}{0.5pt}\end{center}

\section{Complete Example}\label{py-api:complete-example}

\begin{Verbatim}
import stepss

# --- Build test case ---
case = stepss.cfg()
case.addData('dyn_A.dat')
case.addData('settings1.dat')
case.addInit('init.trace')
case.addDst('events.dst')
case.addTrj('output.trj')
case.addObs('obs.dat')
case.addOut('output.trace')
case.addCont('cont.trace')
case.addDisc('disc.trace')

# Runtime observables, recorded by the engine into a curve file
case.addRunObs('BV 1041')
case.addRunObs('MS g1')
case.addRunObs('RT RT')

# Save configuration
case.writeCmdFile('cmd.txt')

# --- Run simulation ---
ram = stepss.sim()

# Start and pause at t = 80 s
ram.execSim(case, 80.0)

# Inject a disturbance dynamically
ram.addDisturb(100.0, 'FAULT BUS 4032 0. 0.')
ram.addDisturb(100.1, 'CLEAR BUS 4032')

# Export Jacobian at this operating point
A, E = ram.getJac()

# Run to end
ram.contSim(ram.getInfTime())

# --- Extract results ---
ext = stepss.extractor(case.getTrj())

# Bus voltages
ext.getBus('4044').mag.plot()

# Generator observables
gen = ext.getSync('g1')
gen.S.plot()    # rotor speed
gen.A.plot()    # rotor angle

# Branch flows
ext.getBranch('1041-01').PF.plot()

# Plot multiple rotor speeds together
stepss.curplot([
    ext.getSync('g1').S,
    ext.getSync('g2').S,
    ext.getSync('g3').S,
])
\end{Verbatim}

\section{Next Steps}\label{py-api:next-steps}

\begin{itemize}
\tightlist
\item
  \hyperref[doc:py-examples]{Examples}, Practical simulation examples and workflows
\item
  \hyperref[doc:ts-index]{Test Systems}, Ready-to-run benchmark systems
\end{itemize}
